\documentclass[11pt]{article}
\newcommand{\SheetCode}{Theory 02}
\usepackage[margin=25mm]{geometry}
\usepackage[T1]{fontenc}
\usepackage{lmodern}
\DeclareUnicodeCharacter{2605}{$\star$}
\usepackage{microtype}
\usepackage{amsmath,amssymb,mathtools,bm}
\usepackage{enumitem}
\usepackage{xcolor}
\usepackage{fancyhdr}
\usepackage{hyperref}
\usepackage[most]{tcolorbox}
\definecolor{agiBlue}{HTML}{173B57}
\definecolor{agiGold}{HTML}{B7791F}
\definecolor{agiLight}{HTML}{EFF6FA}
\definecolor{agiGreen}{HTML}{184D3B}
\hypersetup{colorlinks=true,linkcolor=agiBlue,urlcolor=agiBlue}
\setlength{\parindent}{0pt}
\setlength{\parskip}{0.55em}
\setlength{\headheight}{14pt}
\setlist{nosep,leftmargin=*}
\pagestyle{fancy}
\fancyhf{}
\fancyhead[L]{\small Part of the \href{https://github.com/mmikhasenko/GIModel.jl}{Agentic Gottried--Isgur} project}
\fancyhead[R]{\SheetCode}
\fancyfoot[C]{\thepage}
\newcommand{\dd}{\,\mathrm d}
\newcommand{\ket}[1]{\lvert #1\rangle}
\newcommand{\bra}[1]{\langle #1\rvert}
\newcommand{\braket}[2]{\langle #1\mid #2\rangle}
\newcommand{\expect}[1]{\left\langle #1\right\rangle}
\newcommand{\op}[1]{\widehat{#1}}
\newcommand{\GeV}{\,\mathrm{GeV}}
\newcommand{\R}{\mathbb{R}}
\newtcolorbox{goals}{colback=agiLight,colframe=agiBlue,title=Learning outcomes}
\newtcolorbox{reading}{colback=white,colframe=agiGold,title=Book reference,
  after skip=5mm}
\newtcolorbox{paperbridge}{colback=agiGreen!7,colframe=agiGreen,
  colbacktitle=agiGreen,coltitle=white,title=Connection to the GI paper}
\newtcolorbox{problem}[1]{colback=white,colframe=agiBlue,title={#1},
  before skip=3.5mm,after skip=1mm}
\newtcolorbox{solution}{colback=black!2,colframe=black!45,title=Solution}
\newtcolorbox{quiz}{colback=white,colframe=agiGold,title=Post-solution concept check}
\newtcolorbox{checkpoint}{colback=white,colframe=agiGold,title=Interpretation checkpoint}
\newcommand{\sheettitle}[3]{%
  \begin{center}
  {\Large\bfseries #1}\par
  \vspace{0.2em}{\color{agiBlue}#2}\par
  \vspace{0.4em}\small #3
  \end{center}\vspace{0.5em}}
\newcommand{\difficulty}[1]{\textbf{Difficulty:} #1}
\newcommand{\timebox}[1]{\textbf{Suggested time:} #1}
\newcommand{\answer}{\par\smallskip\color{black!50}\textbf{Answer:} }

\begin{document}
\sheettitle{Sheet 2: Angular momentum and meson identity}{How $n^{2S+1}L_J$ becomes $J^{PC}$}{Prerequisite: Sheet 1. \difficulty{★--★★★} \quad \timebox{2--3 hours}}

\begin{goals}
Couple two spin-$1/2$ constituents to orbital angular momentum; derive meson
parity and charge conjugation; identify which channels may mix; compute the
basic spin expectation values used throughout GI.
\end{goals}
\begin{reading}
Selected reading in Landau--Lifshitz, \emph{Quantum Mechanics: Non-Relativistic
Theory}: Ch. IV, ``Angular Momentum,'' Ch. VIII, ``Spin,'' Ch. IX, ``Identity
of Particles,'' and Ch. XIV, ``Addition of Angular Momenta.'' This is a
reference selection, not the organizing structure of the sheet.
\end{reading}

The spectroscopic notation is $n^{2S+1}L_J$ with $L=S,P,D,F,\ldots$ for
$0,1,2,3,\ldots$.


\begin{problem}{Problem 1 ★ — Enumerate allowed $J$ values}
Prove that coupling orbital $L$ and total spin $S$ permits
$J=|L-S|,\ldots,L+S$. List the $S$- and $P$-wave $q\bar q$ multiplets.
\end{problem}
\begin{solution}
The triangle rule for addition of angular momenta gives the range. For $L=0$:
$^1S_0$ and $^3S_1$. For $L=1$: $^1P_1$ and
$^3P_{0,1,2}$. The radial label $n$ distinguishes repeated multiplets but does
not affect angular quantum numbers.
\end{solution}
\begin{quiz}
Why can $^1P_1$ and $^3P_1$ mix while $^1P_1$ and $^3P_2$ cannot? \answer a
rotational scalar Hamiltonian preserves $J$.
\end{quiz}

\begin{problem}{Problem 2 ★★ — Derive $P$ and $C$}
For a fermion--antifermion state, derive $P=(-1)^{L+1}$. For a neutral
self-conjugate flavor state, let charge conjugation interchange the quark and
antiquark and derive $C=(-1)^{L+S}$ by keeping track of the orbital, spin, and
fermionic exchange signs. Compute $J^{PC}$ for the six multiplets above.
\end{problem}
\begin{solution}
Spatial inversion gives $Y_{Lm}(-\widehat{\bm r})=(-1)^LY_{Lm}$, while the
intrinsic parities of a fermion and its antifermion have product $-1$. Hence
$P=(-1)^{L+1}$.

For charge conjugation, write the spin wavefunction in the tensor-product basis
$\ket{s_1}\otimes\ket{s_2}$. Interchanging the two spin labels gives
\[
P_{12}^{\rm spin}\ket{S M_S}=(-1)^{S+1}\ket{S M_S}:
\]
the $S=1$ triplet is symmetric and the $S=0$ singlet is antisymmetric. The same
interchange sends the relative coordinate $\bm r=\bm r_q-\bm r_{\bar q}$ to
$-\bm r$, contributing $(-1)^L$. Finally, in a state written with creation
operators $b^\dagger d^\dagger\ket0$, charge conjugation first produces
$d^\dagger b^\dagger\ket0=-b^\dagger d^\dagger\ket0$; the last equality is the
fermionic anticommutation sign. Therefore
\[
C=-(-1)^L(-1)^{S+1}=(-1)^{L+S}.
\]
Thus
$^1S_0:0^{-+}$, $^3S_1:1^{--}$, $^1P_1:1^{+-}$, and
$^3P_{0,1,2}:0^{++},1^{++},2^{++}$.
\end{solution}
\begin{quiz}
Why is $C$ not a label for a charged or open-flavor meson? \answer charge
conjugation changes its flavor content, so it is not an eigenstate of $C$.
\end{quiz}

\begin{problem}{Problem 3 ★ — Compute $\bm S_1\cdot\bm S_2$}
Begin with the tensor-product states
$\ket\uparrow\otimes\ket\uparrow$, $\ket\uparrow\otimes\ket\downarrow$,
$\ket\downarrow\otimes\ket\uparrow$, and
$\ket\downarrow\otimes\ket\downarrow$. Construct the three symmetric states
with total spin $S=1$ (the triplet) and the one antisymmetric state with $S=0$
(the singlet). For a normalized state $\ket\psi$, the notation
$\expect O_\psi\equiv\bra\psi O\ket\psi$ means the expectation value, or mean
result of repeated measurements, of $O$. Compute
$\expect{\bm S_1\cdot\bm S_2}_{S=0}$ and
$\expect{\bm S_1\cdot\bm S_2}_{S=1}$.
\end{problem}
\begin{solution}
The coupled spin states are
\[
\begin{aligned}
\ket{1,1}&=\ket{\uparrow\uparrow},&
\ket{1,0}&=\frac{\ket{\uparrow\downarrow}+\ket{\downarrow\uparrow}}{\sqrt2},&
\ket{1,-1}&=\ket{\downarrow\downarrow},\\
\ket{0,0}&=\frac{\ket{\uparrow\downarrow}-\ket{\downarrow\uparrow}}{\sqrt2}.
\end{aligned}
\]
From $\bm S^2=(\bm S_1+\bm S_2)^2$,
\[
\bm S_1\cdot\bm S_2=\frac12[S(S+1)-\tfrac34-\tfrac34].
\]
The result is $-3/4$ for $S=0$ and $+1/4$ for $S=1$.
\end{solution}
\begin{quiz}
Let $E_0$ be the energy before the contact correction and let its positive
radial coefficient be $A>0$, so
$\Delta E=A\expect{\bm S_1\cdot\bm S_2}$. Which corrected energy is higher?
\answer $E_{S=1}=E_0+A/4$ is increased above $E_0$, while
$E_{S=0}=E_0-3A/4$ is decreased below it. Thus ``pushed upward'' means that the
corrected energy is larger than the uncorrected energy.
\end{quiz}

\begin{problem}{Problem 4 ★ — Compute $\bm L\cdot\bm S$}
Derive its diagonal expectation in a state $\ket{LSJM}$ and tabulate the
$^3P_J$ values.
\end{problem}
\begin{solution}
$2\bm L\cdot\bm S=\bm J^2-\bm L^2-\bm S^2$, hence
\[
\expect{\bm L\cdot\bm S}=\tfrac12[J(J+1)-L(L+1)-S(S+1)].
\]
For $L=S=1$, the values at $J=0,1,2$ are $-2,-1,+1$.
\end{solution}
\begin{quiz}
Why is the spin--orbit energy correction zero for an $S$ wave? \answer an
$S$ wave has $L=0$, so $\bm L\ket{L=0,m=0}=0$. Consequently
$\bm L\cdot\bm S$ has zero expectation value, independently of the radial wave.
\end{quiz}

\begin{problem}{Problem 5 ★★ — Prove a centroid sum rule}
For fixed total angular momentum $J$, its projection $M_J$ can take the values
$-J,-J+1,\ldots,J$, giving $2J+1$ states with the same $J$. These are called
magnetic substates. The centroid of a multiplet is the energy average weighted
by this number of states. Show that
$\sum_J(2J+1)\expect{\bm L\cdot\bm S}_J=0$ for fixed $L,S$. Verify it for
$^3P_J$.
\end{problem}
\begin{solution}
Before coupling $\bm L$ and $\bm S$, the relevant Hilbert space has basis
$\ket{Lm_L}\otimes\ket{Sm_S}$; this is the tensor-product space
$\mathcal H_L\otimes\mathcal H_S$. In this basis
\[
\bm L\cdot\bm S=L_zS_z+\tfrac12(L_+S_-+L_-S_+).
\]
The ladder terms change $m_L$ and $m_S$, so their diagonal matrix elements are
zero. Summing the remaining diagonal elements gives
\[
\sum_{m_L,m_S}m_Lm_S=
\left(\sum_{m_L=-L}^{L}m_L\right)
\left(\sum_{m_S=-S}^{S}m_S\right)=0.
\]
The trace --- the sum of diagonal matrix elements --- is unchanged by changing
to the coupled basis $\ket{LSJM}$. In that basis $\bm L\cdot\bm S$ has one
eigenvalue for each $J$, repeated for its $2J+1$ values of $M$. Therefore its
trace is exactly $\sum_J(2J+1)\expect{\bm L\cdot\bm S}_J=0$. For $^3P_J$:
$1(-2)+3(-1)+5(+1)=0$.
\end{solution}
\begin{quiz}
What does a violation of this sum rule in a calculation most directly suggest?
\answer an incorrect angular coefficient, degeneracy factor $(2J+1)$, or state
label. The identity contains no radial integral, so changing the radial
wavefunction cannot repair it.
\end{quiz}

\begin{problem}{Problem 6 ★★★ — Decide which GI mixings are allowed}
Using only $J$, parity, charge conjugation, and flavor, decide whether each pair
can mix: (a) $^1P_1$--$^3P_1$ in $c\bar q$; (b) the same pair in $c\bar c$;
(c) $^3S_1$--$^3D_1$; (d) $n\bar n$--$s\bar s$ pseudoscalars.
\end{problem}
\begin{solution}
(a) Yes: equal $J^P$ and no good $C$; unequal masses allow antisymmetric
spin--orbit coupling. (b) No in a $C$-invariant Hamiltonian: $C$ differs, and
the antisymmetric mass factor vanishes. (c) Yes: tensor forces preserve
$J,S$ and parity while allowing $\Delta L=2$. (d) Yes: the quantum numbers
match and gluon annihilation couples self-conjugate flavor channels.
\end{solution}
\begin{quiz}
The phrase ``same $J$'' is necessary but not sufficient for mixing. Name two
more tests. \answer matching exact discrete symmetries and a nonzero matrix
element for the proposed interaction.
\end{quiz}

\begin{paperbridge}
\small
GI Eq. (4) contains $\bm S_1\cdot\bm S_2$ and its tensor partner. Eq. (14)
diagonalizes in $\ket{jm;ls}$ sectors, with $\bm S=\bm S_1+\bm S_2$ and
$\bm J=\bm L+\bm S$. The paragraph after Eq. (14) identifies the same-$J$
tensor and antisymmetric spin--orbit mixings; Figs. 3--11 use
$n^{2S+1}L_J$ labels.
\end{paperbridge}
\end{document}
